\documentclass[11pt,a4paper]{article}
\usepackage[T1]{fontenc}
\usepackage[utf8]{inputenc}
\usepackage{lmodern}
\usepackage[a4paper,margin=25mm,headheight=15pt]{geometry}
\usepackage{amsmath,amssymb,amsthm,mathtools}
\usepackage{microtype,booktabs,longtable,array}
\usepackage[dvipsnames]{xcolor}
\usepackage{enumitem,fancyhdr}
\usepackage[colorlinks=true,linkcolor=MidnightBlue,citecolor=MidnightBlue,urlcolor=MidnightBlue]{hyperref}
\usepackage{listings}
\definecolor{Ink}{HTML}{15354B}
\definecolor{Soft}{HTML}{F2F5F7}
\hypersetup{pdftitle={Natural densities of the three additive complexities of the Tribonacci word},pdfauthor={AI-assisted research draft},pdfsubject={A proposed resolution of the proportion question in Remark 19 of Popoli--Shallit--Stipulanti}}
\pagestyle{fancy}
\fancyhf{}
\fancyhead[L]{\small\textsc{Tribonacci additive-complexity densities}}
\fancyhead[R]{\small Research draft}
\fancyfoot[C]{\thepage}
\renewcommand{\headrulewidth}{0.3pt}
\setlength{\parskip}{0.35em}
\setlength{\parindent}{1.2em}
\setlength{\emergencystretch}{2em}
\setlist{itemsep=0.15em,topsep=0.4em}
\lstset{basicstyle=\small\ttfamily,columns=fullflexible,keepspaces=true,breaklines=true,frame=single,rulecolor=\color{black!20},backgroundcolor=\color{Soft},showstringspaces=false}
\newtheorem{theorem}{Theorem}[section]
\newtheorem{lemma}[theorem]{Lemma}
\newtheorem{proposition}[theorem]{Proposition}
\newtheorem{corollary}[theorem]{Corollary}
\theoremstyle{definition}
\newtheorem{definition}[theorem]{Definition}
\newtheorem{example}[theorem]{Example}
\theoremstyle{remark}
\newtheorem{remark}[theorem]{Remark}
\newcommand{\N}{\mathbb N}
\newcommand{\Z}{\mathbb Z}
\newcommand{\Q}{\mathbb Q}
\newcommand{\C}{\mathbb C}
\newcommand{\ttt}{\mathbf t}
\newcommand{\one}{\mathbf 1}
\newcommand{\Dec}{\operatorname{Dec}}
\newcommand{\rep}{\operatorname{rep}}
\newcommand{\val}{\operatorname{val}}
\newcommand{\rot}{\operatorname{rot}}
\newcommand{\card}[1]{\lvert #1\rvert}
\newcommand{\code}[1]{\texttt{#1}}
\newcommand{\cyl}{\mathcal B}
\newcommand{\cc}{\mathcal C}
\newcommand{\dd}{\mathrm d}
\DeclareMathOperator{\sgn}{sgn}

\begin{document}
\begin{titlepage}
\centering
\vspace*{12mm}
{\small\scshape Combinatorics on words\par}
\vspace{8mm}
{\LARGE\bfseries\color{Ink} Natural densities of the three\par
additive complexities of the\par Tribonacci word\par}
\vspace{6mm}
{\large A proposed resolution of the proportion question\par
in Popoli--Shallit--Stipulanti, Remark 19\par}
\vspace{8mm}
{AI-assisted research draft\par}
{20 September 2026\par}
\vspace{10mm}
\begin{minipage}{0.92\textwidth}
\begin{abstract}
The Tribonacci word is the fixed point of $0\mapsto01$, $1\mapsto02$,
$2\mapsto0$. Its additive complexity $a(n)$ counts the distinct sums of
length-$n$ factors. Popoli, Shallit and Stipulanti proved that the only values
for positive $n$ are $3,4,5$, and asked for their respective proportions.
We give an exact finite-certificate proof of the existence of the ordinary
natural densities and evaluate them in the cubic field of the Tribonacci
constant $\beta$, where $\beta^3=\beta^2+\beta+1$:
\[
 d_3=\frac{21\beta^2-38\beta+5}{22},\quad
 d_4=\frac{-55\beta^2+72\beta+67}{22},\quad
 d_5=\frac{17\beta^2-17\beta-25}{11}.
\]
For every integer cutoff $N$, not merely for Tribonacci-number cutoffs, the
count of lengths with value $j$ is $d_jN+O(N^\theta)$, where
$\theta=\log(7/5)/\log\beta<0.553$.
The proof reconstructs a 76-state output automaton using abelian
co-decomposition, verifies an exact algebraic left eigenvector and an
integer matrix annihilator, and controls arbitrary cutoffs by a disjoint
prefix-cylinder decomposition. The package includes the full transition and
eigenvector tables, a standard-library-only verifier, exact counting code,
and a separate all-factor check through length $2048$.
\end{abstract}
\vspace{3mm}
\noindent\fcolorbox{black!20}{Soft}{\parbox{0.94\linewidth}{\small
\textbf{Status.} This is a proposed solution with reproduced exact finite
certificates. It has not been independently refereed or formalized in a proof
assistant. The literature audit found the explicit open question and no later
resolution in the sources checked; it is not a guarantee of global priority.
The known range $\{3,4,5\}$ and the co-decomposition method are not claimed as new.}}
\end{minipage}
\vfill
{\small Source, data, reproducibility notes and code accompany this article.\par}
\end{titlepage}

{\small\setlength{\parskip}{0pt}\tableofcontents}
\newpage

\section{The problem and the result}
\label{sec:intro}

\subsection{The invariant being counted}
Let $\varphi$ be the substitution
\begin{equation}\label{eq:phi}
 \varphi(0)=01,\qquad \varphi(1)=02,\qquad \varphi(2)=0,
\end{equation}
and let
\[
 \ttt=\lim_{k\to\infty}\varphi^k(0)
      =010201001020101020100\cdots
\]
be its one-sided fixed point. Positions are indexed from zero. We write
$\ttt[i:i+n]$ for the factor starting at $i$ and having length $n$.
For a finite word $u=u_0\cdots u_{m-1}$ on $\{0,1,2\}$, set
\[
 \sigma(u)=\sum_{i=0}^{m-1}u_i,
 \qquad
 \Psi(u)=(|u|_0,|u|_1,|u|_2).
\]
The additive complexity is
\begin{equation}\label{eq:additive}
 a(n)=\#\{\sigma(\ttt[i:i+n]):i\ge0\}.
\end{equation}
Thus $a(n)$ counts \emph{different sums}, not factors, occurrences, or
Parikh vectors. In particular, two different Parikh vectors can contribute
the same additive class. We use $a(0)=1$ for the empty factor.

For example, the letters all occur, so $a(1)=3$. The length-two factors are
$00,01,02,10,20$; their sums are $0,1,2$, so $a(2)=3$. By contrast, there
are five distinct length-two factors and three distinct length-two Parikh
vectors. These distinctions matter when interpreting a frequency question.

For $j\in\{3,4,5\}$ define
\begin{equation}\label{eq:countdef}
 C_j(N)=\#\{0\le n<N:a(n)=j\}.
\end{equation}
The variable being sampled uniformly is the \emph{length} $n$, not the
starting position of a factor. We seek the limits $C_j(N)/N$.

\subsection{The published question and the exclusion check}
In Theorem~18 of their full preprint \cite{PSS}, Popoli, Shallit and
Stipulanti establish the range $\{3,4,5\}$ for positive lengths and provide
a 76-state Tribonacci automaton computing the function. Their Remark~19
then leaves the respective proportions open. Our question is precisely
that proportion question, with ``proportion'' made explicit as ordinary
natural density over integer lengths.

The area was selected before this question was chosen: a single call to
Python's OS-backed \code{secrets.randbelow(120)} returned zero-based index
$67$, selecting \emph{Combinatorics on words}, item $68$ of the recorded
120-area list. The complete list and result are in
\code{data/selection.json}. The uploaded catalogue
\code{manifest(1).tex}, containing 71 research packages, was read in full.
This target is not among its entries. Its transfinite-word order-type
problems and its discrepancy problem for $0\mapsto1$, $1\mapsto(10)^m$
concern different substitutions, invariants and questions.

A focused current-literature search on 20 September 2026 found no later
resolution of this particular density problem. The checked sources include
the full 2024 preprint and the later work of Couvreur et al.
\cite{Couvreur}, in its April 2025 version~2. The latter develops more general
automata constructions for abelian-type complexities; it is not used as a
proof of the densities below. The scope and exact source versions are
recorded in \code{notes/sources.md}. Failure to find another solution is not
itself evidence sufficient to assert absolute novelty.

\subsection{Main theorem}
Let $\beta$ be the real root greater than one of
\begin{equation}\label{eq:beta}
 f(x)=x^3-x^2-x-1.
\end{equation}
Numerically, $\beta=1.83928675521416113255\ldots$.

\begin{theorem}[Natural densities and a power-saving error]\label{thm:main}
For the additive complexity \eqref{eq:additive}, the ordinary natural
densities of the values $3,4,5$ exist and equal
\begin{align}
 d_3&=\frac{21\beta^2-38\beta+5}{22},\label{eq:d3}\\
 d_4&=\frac{-55\beta^2+72\beta+67}{22},\label{eq:d4}\\
 d_5&=\frac{17\beta^2-17\beta-25}{11}.\label{eq:d5}
\end{align}
More precisely, with
\begin{equation}\label{eq:theta}
 \theta=\frac{\log(7/5)}{\log\beta}
       =0.55215697321855134306\ldots,
\end{equation}
we have, for each $j\in\{3,4,5\}$,
\begin{equation}\label{eq:mainerror}
 C_j(N)=d_jN+O(N^\theta)\qquad(N\to\infty).
\end{equation}
All three densities are strictly positive and their sum is one.
\end{theorem}

The decimal values are
\begin{center}
\begin{tabular}{@{}crr@{}}
\toprule
Value & Natural density & Percentage\\
\midrule
$3$ & $0.27952701944967565180\ldots$ & $27.9527019449\ldots\%$\\
$4$ & $0.60749905184438815304\ldots$ & $60.7499051844\ldots\%$\\
$5$ & $0.11297392870593619515\ldots$ & $11.2973928705\ldots\%$\\
\bottomrule
\end{tabular}
\end{center}

The main point beyond extracting a dominant eigenvector is that
\eqref{eq:mainerror} holds at \emph{every} cutoff. Limits along a sparse
sequence of numeration-system cutoffs do not by themselves establish natural
density. Section~\ref{sec:allcutoffs} supplies the missing uniform step.

\subsection{What is new, what is reused, and what is certified}
The range of $a(n)$ and the existence of a 76-state automaton are already
known \cite{PSS}. Abelian co-decomposition is a known mechanism, developed
for this setting by Turek \cite{Turek}. We specialize and prove the needed
recursions, then reconstruct the additive-complexity machine independently.
We do not claim that our state numbering or serialized machine is identical
to the authors' original file.

The proposed contribution is the exact evaluation of the three densities,
the uniform all-cutoff proof, and the explicit power-saving estimate.
The finite parts of the argument are exposed as reproducible certificates:
a closed family of co-decompositions, the quotient map to the displayed
machine, an algebraic eigenvector, and a polynomial identity for the integer
adjacency matrix. A verifier checks those finite identities using only
Python's standard library. The passage from them to an infinite-word theorem
and a limit theorem is proved in the text.

\section{Tribonacci numeration and word prefixes}
\label{sec:numeration}

Define words and lengths
\[
 W_k=\varphi^k(0),\qquad T_k=|W_k|.
\]
The initial values and recursions are
\begin{equation}\label{eq:Trec}
 T_0=1,\quad T_1=2,\quad T_2=4,\qquad
 T_k=T_{k-1}+T_{k-2}+T_{k-3}\quad(k\ge3),
\end{equation}
and
\begin{equation}\label{eq:Wrec}
 W_k=W_{k-1}W_{k-2}W_{k-3}\quad(k\ge3).
\end{equation}
Indeed, $\varphi^k(0)=\varphi^{k-1}(0)\varphi^{k-1}(1)$,
$\varphi^{k-1}(1)=\varphi^{k-2}(0)\varphi^{k-2}(2)$, and
$\varphi^{k-2}(2)=\varphi^{k-3}(0)$.

\begin{lemma}[Canonical representations]\label{lem:numeration}
For every $k\ge0$, the binary strings of length $k$ with no occurrence of
$111$ represent the integers $0\le n<T_k$ exactly once under
\[
 \val(b_{k-1}\cdots b_0)=\sum_{i=0}^{k-1}b_iT_i.
\]
Their lexicographic order agrees with numerical order. Leading zeros do not
change the represented number. Every positive integer therefore has a unique
representation beginning with $1$ and avoiding $111$.
\end{lemma}
\begin{proof}
The assertions are immediate for $k=0,1,2$. For $k\ge3$, the allowed strings
split into the three disjoint prefix classes $0$, $10$, and $110$.
Their suffixes have lengths $k-1,k-2,k-3$, respectively, and may be arbitrary
allowed strings of those lengths. By induction their represented intervals
are
\[
 [0,T_{k-1}),\quad
 [T_{k-1},T_{k-1}+T_{k-2}),\quad
 [T_{k-1}+T_{k-2},T_k).
\]
These intervals are adjacent, disjoint and ordered exactly as the three
prefixes. This proves all assertions, including uniqueness and order.
\end{proof}

Write $P_n=\ttt[0:n]$ and use $\rep(n)$ for the canonical representation,
with $\rep(0)$ the empty string.

\begin{lemma}[The prefix associated with a representation]\label{lem:prefix}
If $\rep(n)=b_s\cdots b_0$, then
\begin{equation}\label{eq:prefix}
 P_n=W_s^{b_s}W_{s-1}^{b_{s-1}}\cdots W_0^{b_0}.
\end{equation}
Consequently, if appending $d\in\{0,1\}$ gives another allowed
representation, and $N=\val(\rep(n)d)$, then
\begin{equation}\label{eq:prefixtransition}
 P_N=\varphi(P_n)0^d.
\end{equation}
\end{lemma}
\begin{proof}
The base cases follow from $W_0=0$, $W_1=01$, $W_2=0102$.
For a prefix shorter than $T_k$, the same three intervals in the proof of
Lemma~\ref{lem:numeration} say that it lies in the first block of
\eqref{eq:Wrec}, consists of that block and a prefix of the second, or
consists of the first two and a prefix of the third. Induction within those
blocks gives \eqref{eq:prefix}. Applying $\varphi$ replaces each $W_i$
by $W_{i+1}$; an appended digit $d$ contributes $W_0^d=0^d$.
This proves \eqref{eq:prefixtransition}.
\end{proof}

The generating function of the numeration weights is
\begin{equation}\label{eq:Tgf}
 \sum_{k\ge0}T_kz^k=\frac{1+z+z^2}{1-z-z^2-z^3}.
\end{equation}
For later use it is enough that $T_k\asymp\beta^k$. This also follows below
from the exact automaton spectrum; no asymptotic estimate is being assumed
without proof.

\section{From all factors to a finite automaton}
\label{sec:automaton}

\subsection{Co-decomposition and relative sums}
\begin{definition}
For finite words $u,v$ with $\Psi(u)=\Psi(v)$, cut both words at every
position $i$ at which $\Psi(u[0:i])=\Psi(v[0:i])$. The set of paired
nonempty pieces so obtained is denoted $\Dec(u,v)$. Repeated identical
pairs are retained only once. This is the finest, or maximal-cut,
abelian co-decomposition.
\end{definition}
The use of a set is intentional: only which relative vectors occur matters,
not how many times a particular piece appears.

For a pair $(u,v)$ of equal length define its scalar prefix differences by
\[
 S(u,v)=\{\sigma(v[0:i])-\sigma(u[0:i]):0\le i\le|u|\},
\]
and for a finite set $Z$ of such pairs let
\begin{equation}\label{eq:out}
 S(Z)=\bigcup_{(u,v)\in Z}S(u,v),\qquad o(Z)=|S(Z)|.
\end{equation}
If a pair is cut at a common Parikh prefix, the relative vector, and hence
the scalar difference, is zero at the cut. It follows that $S(u,v)$ is
exactly the union of the scalar-difference sets of its pieces.

\subsection{A finite word containing every relevant factor}
Fix $n\ge1$, choose $K$ such that $n\le T_K$, and put
\[
 W=\varphi^{K+3}(0),\qquad
 \rot_n(W)=P_n^{-1}WP_n.
\]
Here $P_n^{-1}W$ means deletion of the specified prefix, not a group
inverse. The prefix is available because $P_n$ is a prefix of $W$.

\begin{lemma}[Coverage and independence]\label{lem:coverage}
The set
\begin{equation}\label{eq:Zn}
 Z(n)=\Dec(W,\rot_n(W))
\end{equation}
is independent of the choice of $K$ satisfying $n\le T_K$. Moreover,
\begin{equation}\label{eq:Zoutput}
 a(n)=o(Z(n)).
\end{equation}
\end{lemma}
\begin{proof}
Set $B_a=\varphi^{K+1}(a)$ for $a\in\{0,1,2\}$. Each $B_a$ begins with
$W_K$: explicitly, $B_2=W_K$, while $B_0$ and $B_1$ begin with $W_K$.
Thus each starts with $P_n$ and has length at least $n$.
The fixed point is a concatenation of these blocks. A length-$n$ factor
starting inside $B_a$ is contained in $B_aP_n$, because any part extending
beyond the block is a prefix of the next block of length at most $n$.

Now $W=\varphi^{K+1}(0102)$ contains every block type $B_a$. Every
length-$n$ factor is therefore found at a position $0\le i<|W|$ in
$WP_n$. Also $WP_n$ really is a prefix of $\ttt$: the letter following the
first $0$ at the substitution level is $1$, and
$\varphi^{K+3}(1)$ begins with $W_{K+2}$, which is long enough to contain
$P_n$.

For $0\le i\le|W|$, subtraction of prefix vectors gives
\begin{align*}
 \Psi(\rot_n(W)[0:i])-\Psi(W[0:i])
 &=\Psi(\ttt[n:n+i])-\Psi(\ttt[0:i])\\
 &=\Psi(\ttt[i:i+n])-\Psi(P_n).
\end{align*}
As $i$ varies, these are exactly the relative Parikh vectors of length-$n$
factors; the possible endpoint $i=|W|$ contributes the already present zero
vector. Cutting at every zero vector does not change the union of the
prefix-difference vectors. Applying the linear functional
$(x_0,x_1,x_2)\mapsto x_1+2x_2$ and adding the constant $\sigma(P_n)$
therefore proves \eqref{eq:Zoutput}.

For independence, let $K_0$ be the least index with $n\le T_{K_0}$ and
instead fix $B_a=\varphi^{K_0+1}(a)$. For every $K\ge K_0$ we have
\[
 \varphi^{K+3}(0)
 =\varphi^{K_0+1}\bigl(\varphi^{K-K_0+2}(0)\bigr).
\]
The word inside the parentheses contains all three letters, since
$\varphi^2(0)=0102$. Each $B_a$ starts with $P_n$. If $B_a=P_nU_a$, then
rotating a concatenation of these blocks by $P_n$ gives the concatenation
of the corresponding $U_aP_n$. At every block boundary the top and bottom
have the same Parikh vector. The finest co-decomposition is consequently
\[
 \bigcup_{a\in\{0,1,2\}}\Dec(B_a,P_n^{-1}B_aP_n),
\]
which no longer depends on $K$. This proves independence.
\end{proof}

\subsection{The digit transition}
The following specializes the co-decomposition mechanism of Turek
\cite{Turek}; we give the argument because it is the connection between the
infinite word and the finite certificate.

For a pair $(u,v)$ define
\begin{align}
 D_0(u,v)&=\Dec(\varphi(u),\varphi(v)),\label{eq:D0}\\
 D_1(u,v)&=\Dec(\varphi(u),0^{-1}\varphi(v)0).\label{eq:D1}
\end{align}
Every nonempty image under $\varphi$ begins with $0$, so the prefix deletion
in \eqref{eq:D1} is always legal. Extend the operation to a set by
$D_d(Z)=\bigcup_{(u,v)\in Z}D_d(u,v)$.

\begin{proposition}[Exact finite-state update]\label{prop:transition}
We have
\begin{equation}\label{eq:initialZ}
 Z(1)=\{(0,0),(01,10),(02,20)\}.
\end{equation}
If $\rep(N)=\rep(n)d$ is allowed and $n\ge1$, then
\begin{equation}\label{eq:Ztransition}
 Z(N)=D_d(Z(n)).
\end{equation}
\end{proposition}
\begin{proof}
The initial calculation uses $W=\varphi^3(0)=0102010$ and its one-letter
rotation $1020100$; cutting at all equal Parikh prefixes gives
\eqref{eq:initialZ}.

Choose $K$ sufficiently large that both the old and the new prefix fit in
the block constructions of Lemma~\ref{lem:coverage}. By
Lemma~\ref{lem:prefix}, $P_N=\varphi(P_n)0^d$. For $d=0$, applying
$\varphi$ to the old paired words gives the new rotated pair. All former
piece boundaries remain common Parikh boundaries, so refinement inside each
image pair gives exactly $D_0(Z(n))$.

For $d=1$, the new bottom word is the additional one-letter rotation of the
bottom word just obtained. If its consecutive image pieces are
$\varphi(v_1),\ldots,\varphi(v_m)$, each starts with $0$. Writing each as
$0V_i$ shows
\[
 0^{-1}\bigl((0V_1)\cdots(0V_m)\bigr)0
   =(V_1 0)\cdots(V_m 0).
\]
Thus the one-letter rotation can be performed in each piece separately.
Each rotated piece still has the same Parikh vector as its top partner.
Refinement at all common Parikh prefixes gives $D_1(Z(n))$, as required.
\end{proof}

\subsection{Exhaustive closure and the 76-state quotient}
Starting with \eqref{eq:initialZ}, apply both transitions to every newly
encountered set until no new set remains. The executed construction reaches
exactly $277$ distinct sets, using $56$ distinct word pairs; the longest
word in a pair has length $24$. These are a finite closed family, not a
sample of the first few lengths. The serialized family, its two outgoing
transitions at every state, and its scalar output sets are all in
\code{data/co\_decomposition.json}.

The raw closure is larger than necessary because it also allows digit
strings containing $111$. To restrict to the canonical language, form the
product with the trailing-one counter $r\in\{0,1,2\}$. A zero resets $r$;
a one increments it, and a one is forbidden when $r=2$. Beginning after the
first digit $1$ gives $296$ reachable product states. Add a zero/empty
initial state of output $1$ and an invalid sink of output $0$.

Partition states by output and repeatedly refine using the signature
\[
 \bigl(\text{output}(q),\,
       \text{class}(\delta(q,0)),\,
       \text{class}(\delta(q,1))\bigr).
\]
After stabilization, the quotient has $77$ states including the invalid
sink. Omitting the sink leaves $76$ live states. Canonical breadth-first
numbering, visiting digit $0$ before digit $1$, gives the table in
Appendix~\ref{app:transitions}. Its initial state is $0$, and a dash denotes
the omitted invalid transition.

\begin{proposition}[Certified machine]\label{prop:machine}
For every $n\ge0$, reading $\rep(n)$ from state $0$ in the table of
Appendix~\ref{app:transitions} produces the output $a(n)$. Leading zeros
may be inserted arbitrarily. In particular $a(n)\in\{3,4,5\}$ for $n>0$.
\end{proposition}
\begin{proof}
The finite closure certificate checks \eqref{eq:D0}--\eqref{eq:D1} for all
$277$ sets. Proposition~\ref{prop:transition} and induction on the number of
digits show that its output is $a(n)$ on every allowed positive
representation. The language product forbids exactly $111$. At every
refinement step, states in one stable class have equal outputs and their
transitions go to equal classes. The included quotient map checks this
commuting property for \emph{every} product state and both digits, not only
for representatives. Thus quotienting preserves all outputs on all finite
inputs. The added initial state handles the empty string and leading
zeros. Finally, all positive-representation states have output $3,4$, or
$5$, as the table shows.
\end{proof}

This proposition reproves the already known range theorem. It is included
to make the spectral calculation independent of an unavailable external
automaton file. The implementation is short enough to audit: it uses only
word substitution, integer prefix differences, finite sets, a queue, and
partition refinement. No search through an unbounded sequence is used to
guess a transition.

\section{An exact spectral certificate}
\label{sec:spectrum}

\subsection{The adjacency matrix and suffix counts}
Let $Q=\{0,1,\ldots,75\}$ be the live state set, let $o(q)$ denote its
output, and define the nonnegative integer matrix
\begin{equation}\label{eq:A}
 A_{qt}=\#\{d\in\{0,1\}:\delta(q,d)=t\}.
\end{equation}
An illegal transition contributes nothing. If both digits have the same
target, their contributions are counted with multiplicity. Row vectors act
on the left. For each $j\in\{1,3,4,5\}$, let $h_j$ be the column vector
whose $q$th entry is $1$ when $o(q)=j$ and $0$ otherwise. Then
\begin{equation}\label{eq:suffixcounts}
 e_qA^kh_j
\end{equation}
counts the legal suffixes of length $k$ that, starting at $q$, finish at
output $j$. In particular,
\begin{equation}\label{eq:cutoffcounts}
 C_j(T_k)=e_0A^kh_j,\qquad T_k=e_0A^k\one
\end{equation}
for $j=3,4,5$.

The key task is to understand \eqref{eq:suffixcounts} \emph{uniformly in
$q$}. There are only $76$ states, but their leading amplitudes are not all
the same. Ignoring that dependence would make the all-cutoff argument
invalid.

\subsection{The closed component and a right eigenvector}
The transition table has the closed strongly connected component
\begin{equation}\label{eq:C}
 \cc=\{5,9,10,11\}\cup\{14,15,\ldots,75\}.
\end{equation}
It contains $66$ states, and every live state can reach it. State $62$ has
a self-loop on digit $0$. These assertions are checked by forward and
reverse graph searches in the verifier. Within $\cc$, there are $10$ states
of output $3$, $41$ of output $4$, and $15$ of output $5$.
\emph{The desired proportions are not these counts divided by $66$.}
Different states are reached with different asymptotic weights.

Each live state has a well-defined trailing-one context $r(q)\in\{0,1,2\}$.
It can be recovered directly from the table: $r(q)=2$ if its $1$-transition
is illegal; otherwise $r(q)=1$ if the target's $1$-transition is illegal;
and otherwise $r(q)=0$. The table verifies
\[
 r(\delta(q,0))=0,
 \qquad r(\delta(q,1))=r(q)+1
\]
whenever the latter transition is legal. Define
\begin{equation}\label{eq:right}
 v_q=
 \begin{cases}
 1,&r(q)=0,\\
 \beta-1,&r(q)=1,\\
 \beta^{-1}=\beta^2-\beta-1,&r(q)=2.
 \end{cases}
\end{equation}
All entries are positive, and $v_0=1$.

\begin{lemma}\label{lem:right}
The vector \eqref{eq:right} satisfies $Av=\beta v$.
\end{lemma}
\begin{proof}
At context $0$, the two target values sum to $1+(\beta-1)=\beta$.
At context $1$, they sum to $1+\beta^{-1}=\beta^2-\beta
=\beta(\beta-1)$. At context $2$, only the zero-transition remains,
with value $1=\beta\beta^{-1}$. These are exactly the three eigenvector
equations. The identity $1+\beta^{-1}=\beta^2-\beta$ follows from
\eqref{eq:beta}.
\end{proof}

\subsection{A left eigenvector in the cubic field}
Appendix~\ref{app:eigenvector} supplies integer triples $(u_q,w_q,z_q)$
with
\begin{equation}\label{eq:left}
 \ell_q=\frac{u_q+w_q\beta+z_q\beta^2}{44}.
\end{equation}
The ten entries outside $\cc$ are zero. The exact finite identities are
\begin{equation}\label{eq:leftcert}
 \ell A=\beta\ell,\qquad \ell\one=1.
\end{equation}
No floating-point calculation is needed to verify them. Reduction modulo
$f(\beta)=0$ sends multiplication by $\beta$ to the integer operation
\begin{equation}\label{eq:triplerule}
 (u,w,z)\longmapsto(z,u+z,w+z).
\end{equation}
For each column $q$ of $A$, add the triples at its incoming sources, with
edge multiplicities, and compare with \eqref{eq:triplerule}. This is a
collection of $228$ integer component equalities. Normalization amounts to
\[
 \sum_q(u_q,w_q,z_q)=(44,0,0).
\]
For a concrete example, the incoming sources for state $62$ are $56,62,71$.
Its triple is $(14,-48,22)$, and the sum of the three incoming triples is
$(22,36,-26)$, exactly its image under \eqref{eq:triplerule}.

The sums over output classes are particularly simple:
\begin{center}
\begin{tabular}{@{}crrr@{}}
\toprule
Output & $\sum u_q$ & $\sum w_q$ & $\sum z_q$\\
\midrule
$1$ & $0$ & $0$ & $0$\\
$3$ & $10$ & $-76$ & $42$\\
$4$ & $134$ & $144$ & $-110$\\
$5$ & $-100$ & $-68$ & $68$\\
\bottomrule
\end{tabular}
\end{center}
Consequently
\begin{equation}\label{eq:leftdensity}
 \ell h_j=d_j\quad(j=3,4,5)
\end{equation}
with the expressions in Theorem~\ref{thm:main}. Direct multiplication of
the left and right vectors, reduced using \eqref{eq:beta}, also gives
\begin{equation}\label{eq:normalization}
 \ell v=5\beta^2-6\beta-5,
 \qquad
 \frac1{\ell v}=\frac{3\beta^2+7\beta+2}{22}>0.
\end{equation}
The second identity can be verified by multiplying its two sides by
$5\beta^2-6\beta-5$ and reducing modulo $f$.

\subsection{An integer annihilating polynomial}
The next certificate controls every eigenvalue, including transient states.
Define
\begin{align}
 p(x)={}&x^{45}(x-1)^3(x^2+x+1)(x^8-x^4+1)
 \notag\\[-1mm]
 &\quad\cdot(x^3-x^2-x-1)(x^3+x^2+x-1)
 \notag\\
 &\quad\cdot(x^{12}-2x^9-2x^5-2x+1).
 \label{eq:p}
\end{align}
It has degree $76$.

\begin{lemma}[Finite annihilator certificate]\label{lem:annihilator}
The matrix in \eqref{eq:A} satisfies
\begin{equation}\label{eq:pA}
 p(A)=0.
\end{equation}
\end{lemma}
\begin{proof}[Exact verification]
Expand the product \eqref{eq:p} by integer convolution, then evaluate it at
$A$ by Horner's rule. Each left multiplication by $A$ is only the addition
of the rows corresponding to the one or two legal targets in the transition
table. Every one of the $76^2=5776$ final entries is zero.
The routine \code{check\_spectral} in \code{code/verify.py} performs these
operations verbatim, with unbounded integers. Thus \eqref{eq:pA} is a
finite explicit integer identity, not a numerical eigenvalue estimate.
An optional symbolic derivation also identifies $p$ as the characteristic
polynomial, but that stronger identification is unnecessary for this proof.
\end{proof}

\begin{lemma}[Separation from the dominant root]\label{lem:gap}
Every zero of $p$ other than $\beta$ has modulus strictly less than $7/5$.
The zero $\beta$ has multiplicity one in $p$.
\end{lemma}
\begin{proof}
The zero factors at $0$ and the roots of unity are harmless. The roots of
$x^8-x^4+1$ lie on the unit circle because
\[
 (x^4+1)(x^8-x^4+1)=x^{12}+1.
\]
For $f(x)=x^3-x^2-x-1$, its derivative is $(3x+1)(x-1)$, and its local
maximum, at $-1/3$, is $-22/27<0$. Hence $f$ has just one real root.
Direct substitution places this root in $(9/5,19/10)$. The other roots
are a conjugate pair whose product is $1/\beta$, so each has modulus
$\beta^{-1/2}<1$. The roots of $x^3+x^2+x-1$ are the reciprocals of the
roots of $f$. Their moduli are $\beta^{-1}$ and $\sqrt\beta$, both less
than $7/5$, since $\beta<19/10<49/25$.

It remains to treat
\[
 R(x)=x^{12}-2x^9-2x^5-2x+1.
\]
At $r=7/5$, the following strictly positive margin is exact:
\begin{equation}\label{eq:radiusmargin}
 r^{12}-2r^9-2r^5-2r-1
 =\frac{199057326}{244140625}>0.
\end{equation}
If $R(z)=0$ and $|z|\ge r$, the triangle inequality would give
\[
 1\le 2|z|^{-3}+2|z|^{-7}+2|z|^{-11}+|z|^{-12}
 \le 2r^{-3}+2r^{-7}+2r^{-11}+r^{-12}<1,
\]
a contradiction. This proves the modulus bound for the last factor.
All other factors are nonzero at $\beta$, and the real root of $f$ is
simple, proving the multiplicity assertion.
\end{proof}

\subsection{The rank-one asymptotic and uniform cancellation}
\begin{lemma}[Simple dominant eigenvalue]\label{lem:simple}
The eigenspace of $A$ at $\beta$ is one-dimensional, spanned by $v$.
Its algebraic multiplicity is also one.
\end{lemma}
\begin{proof}
Given an eigenvector $x$, divide its coordinates by the positive vector
$v$, writing $y_q=x_q/v_q$. The equation $Ax=\beta x$ becomes
\[
 y_q=\sum_t P_{qt}y_t,
 \qquad
 P_{qt}=\frac{A_{qt}v_t}{\beta v_q}.
\]
By Lemma~\ref{lem:right}, every row of $P$ sums to one. On the closed
strongly connected set $\cc$, a real harmonic function is constant: at a
maximum, every successor with positive transition probability must have
that same value, and connectivity propagates it throughout $\cc$.
For a complex harmonic function, apply the argument to real and imaginary
parts. Every live state has a path into $\cc$. Applying the same maximum
and minimum argument to all states shows that a real harmonic function
cannot exceed or fall below the constant already forced on $\cc$.
Thus $y$ is constant everywhere and $x$ is a multiple of $v$.

Finally, \eqref{eq:pA} and the simple factor of $p$ at $\beta$ imply that
no Jordan block at $\beta$ has size greater than one. Together with the
one-dimensional eigenspace, this proves algebraic multiplicity one.
\end{proof}

\begin{proposition}[Uniform suffix estimate]\label{prop:spectral}
In any fixed matrix norm,
\begin{equation}\label{eq:spectral}
 A^k=\beta^k\frac{v\ell}{\ell v}
      +O\bigl((7/5)^k\bigr).
\end{equation}
For every $j\in\{3,4,5\}$, uniformly over all live states $q$,
\begin{equation}\label{eq:cancellation}
 e_qA^k(h_j-d_j\one)=O\bigl((7/5)^k\bigr).
\end{equation}
\end{proposition}
\begin{proof}
Lemma~\ref{lem:simple} and \eqref{eq:leftcert} identify the spectral
projector at $\beta$ as $v\ell/(\ell v)$; the denominator is nonzero by
\eqref{eq:normalization}. All remaining eigenvalues lie strictly inside
the circle of radius $7/5$, by Lemmas~\ref{lem:annihilator}
and~\ref{lem:gap}. In Jordan form their powers are a finite sum of terms
bounded by a polynomial in $k$ times $|\lambda|^k$. For each such
$|\lambda|<7/5$, this is $O((7/5)^k)$. Nilpotent blocks vanish after a
finite number of steps. This proves \eqref{eq:spectral}.

Since $\ell\one=1$ and $\ell h_j=d_j$, the dominant rank-one term
annihilates $h_j-d_j\one$. This gives \eqref{eq:cancellation}. A single
constant works for all states because $Q$ is finite.
\end{proof}

There is no assumption here that the automaton graph itself is one
irreducible component. The transient states have been included throughout,
which is essential when a counting block starts near the initial state.

\begin{corollary}\label{cor:Tasymptotic}
With $c_\beta=(3\beta^2+7\beta+2)/22$, we have
\[
 T_k=c_\beta\beta^k+O((7/5)^k),\qquad
 C_j(T_k)=d_jc_\beta\beta^k+O((7/5)^k).
\]
Also $d_j>0$ for $j=3,4,5$ and $d_3+d_4+d_5=1$.
\end{corollary}
\begin{proof}
Apply \eqref{eq:spectral} to \eqref{eq:cutoffcounts}, using $v_0=1$ and
\eqref{eq:normalization}. The sum of the three displayed formulas for $d_j$
is exactly one. Positivity can be seen without relying on decimal
approximations: the matrices $\beta^{-k}A^k$ are nonnegative and converge
to $v\ell/(\ell v)$. Since $v$ and $\ell v$ are positive, $\ell$ is
nonnegative. Its sum is one, and it is supported on $\cc$. The left
eigenvector equation and strong connectivity force it to be strictly
positive at every state of $\cc$. Each of the three output classes meets
$\cc$, so every $\ell h_j$ is positive.
\end{proof}

\section{From numeration cutoffs to every integer cutoff}
\label{sec:allcutoffs}

The passage in this section is what turns asymptotic path counts into the
natural-density theorem. It also explains why a numerical fit along $T_k$
would have been insufficient.

\subsection{The lexicographic decomposition}
Take $N\ge1$ and write its canonical representation as a length-$L$ word
\[
 b=b_{L-1}\cdots b_0.
\]
Pad every representation of an integer below $N$ with leading zeros to
length $L$. For each position at which $b_k=1$, consider the set of words
which agree with $b$ at all more significant positions, have $0$ at
position $k$, and then have any allowed suffix of length $k$. Call this
set $\cyl_k$.

\begin{lemma}[Disjoint prefix cylinders]\label{lem:cylinders}
The allowed length-$L$ words representing $0\le n<N$ are the disjoint union
of the $\cyl_k$ over positions with $b_k=1$. If $q_k$ is the automaton state
after the retained prefix followed by that $0$, then
\[
 |\cyl_k|=e_{q_k}A^k\one,
 \qquad
 \#\{n\in\cyl_k:a(n)=j\}=e_{q_k}A^kh_j.
\]
\end{lemma}
\begin{proof}
Every smaller word has a unique most significant position where it first
differs from $b$. At that position $b$ has $1$ and the smaller word has
$0$. Conversely, any word formed in this way is lexicographically smaller,
and Lemma~\ref{lem:numeration} says it is numerically smaller. The initial
prefix is allowed because it is a prefix of $b$ with a $1$ replaced at its
end by $0$; that replacement cannot create $111$. Legal completions are
precisely the suffix paths counted by $A^k$. Uniqueness of the first
position proves disjointness.
\end{proof}

\begin{proof}[Proof of Theorem~\ref{thm:main}]
Subtract $d_j$ times the first count of Lemma~\ref{lem:cylinders} from the
second and sum over the disjoint cylinders. Proposition~\ref{prop:spectral}
gives
\begin{align*}
 |C_j(N)-d_jN|
 &\le K\sum_{\substack{0\le k<L\\b_k=1}}(7/5)^k\\
 &\le K\sum_{k=0}^{L-1}(7/5)^k
 =O((7/5)^L),
\end{align*}
where $K$ is independent of the prefix state and of $N$.
There is at most one block of each suffix length, so the geometric sum
does not introduce an extra logarithmic factor.

By Lemma~\ref{lem:numeration},
$T_{L-1}\le N<T_L$. Corollary~\ref{cor:Tasymptotic} implies
$\beta^L\asymp N$. With $\theta$ from \eqref{eq:theta},
\[
 (7/5)^L=(\beta^L)^\theta=O(N^\theta).
\]
This proves \eqref{eq:mainerror}. Since $\theta<1$, division by $N$
proves the ordinary natural-density limits. Positivity and normalization
were proved in Corollary~\ref{cor:Tasymptotic}.
\end{proof}

\subsection{Two uniform consequences}
\begin{corollary}[Fixed-prefix proportions]\label{cor:prefix}
Fix any allowed prefix $w$, with leading zeros permitted. Among all allowed
words $ws$ with $|s|=k$, the proportion whose represented length has
additive complexity $j$ tends to $d_j$ as $k\to\infty$.
\end{corollary}
\begin{proof}
If $q$ is the state reached by $w$, the numerator is $e_qA^kh_j$ and the
denominator is $e_qA^k\one$. Their leading terms from
\eqref{eq:spectral} are, respectively,
$\beta^kv_qd_j/(\ell v)$ and $\beta^kv_q/(\ell v)$.
The latter is positive and the error is of smaller exponential order.
\end{proof}

\begin{corollary}[Intervals and the mean complexity]\label{cor:mean}
For integers $X\ge1$ and $1\le H\le X$,
\[
 \#\{X\le n<X+H:a(n)=j\}=d_jH+O(X^\theta).
\]
In particular the corresponding proportions tend to $d_j$ whenever
$H/X^\theta\to\infty$. Furthermore,
\begin{equation}\label{eq:mean}
 \frac1N\sum_{0\le n<N}a(n)
 =\frac{13\beta^2+4\beta+33}{22}+O(N^{\theta-1}),
\end{equation}
whose limiting value is $3.83344690925626054335\ldots$.
\end{corollary}
\begin{proof}
Subtract \eqref{eq:mainerror} at $X$ from the same formula at $X+H$;
because $H\le X$, both error terms are $O(X^\theta)$.
For the mean, the length-zero contribution is $1$, and every positive
length contributes one of $3,4,5$. Thus
\[
 \sum_{n<N}a(n)=1+3C_3(N)+4C_4(N)+5C_5(N)
\]
for $N\ge1$. Substitute the three count estimates and simplify
$3d_3+4d_4+5d_5$ in the basis $1,\beta,\beta^2$.
\end{proof}

\begin{remark}[What the exponent does not say]
The exponent in \eqref{eq:theta} is a certified convenient bound, not an
assertion of optimality. It comes from the rational spectral radius bound
$7/5$, whose proof is the single exact inequality
\eqref{eq:radiusmargin}. The theorem supplies no claimed numerical value
for the implicit constant in the error term, and it does not assert uniform
relative accuracy on intervals shorter than the scale covered by
Corollary~\ref{cor:mean}.
\end{remark}

\section{Exact counting and rational generating functions}
\label{sec:computation}

\subsection{An exact algorithm for any cutoff}
The proof gives a practical counting algorithm. Let
\[
 F_k(q,j)=e_qA^kh_j.
\]
Its initialization and recurrence are
\begin{equation}\label{eq:DP}
 F_0(q,j)=\mathbf 1_{o(q)=j},\qquad
 F_{k+1}(q,j)=\sum_{\substack{d\in\{0,1\}\\\delta(q,d)\text{ legal}}}
 F_k(\delta(q,d),j).
\end{equation}
Precompute these arrays for $0\le k\le L$, read the digits of $N$, and at
each digit $1$ add the array for the alternative zero-transition and the
remaining suffix length. This is exactly the disjoint union used in
Lemma~\ref{lem:cylinders}. Only additions of nonnegative integers are
required; no density approximation enters the computation.

There are $76$ states and four output labels including the empty-factor
value $1$. The precomputation uses $O(76L)$ integer additions up to a fixed
constant, and the final digit scan uses $O(L)$ more. The integers have
$O(L)$ bits. With ordinary addition this gives a conservative
$O(76L^2)$ bit-operation bound, where $L=O(\log N)$; retaining all suffix
layers uses $O(76L^2)$ bits. Computing a single $a(n)$ uses one transition
per digit after conversion to Tribonacci numeration.

Here are exact counts produced by the included code. The three displayed
counts sum to $N-1$, because the single length $0$ has value $1$.
\begin{center}
\begin{tabular}{@{}rrrr@{}}
\toprule
$N$ & $C_3(N)$ & $C_4(N)$ & $C_5(N)$\\
\midrule
\csname @@input\endcsname tables/counts_rows.tex 
\bottomrule
\end{tabular}
\end{center}
The data files also include exact counts at $10^{18}$, $10^{30}$ and
$10^{100}$. These are finite computations of whole intervals, not
extrapolations from smaller samples. For example, at $N=10^{18}$ the counts
are
\begin{align*}
 C_3(10^{18})&=279527022667110675,\\
 C_4(10^{18})&=607499052340780247,\\
 C_5(10^{18})&=112973924992109077.
\end{align*}
The comparatively slow visible convergence at moderate cutoffs is compatible
with the error bound; it was not used to guess exact coefficients in the
density formulas.

\subsection{Generating functions at Tribonacci cutoffs}
The whole sequence of cutoff counts has a compact rational description.
Set
\begin{equation}\label{eq:cutoffGF}
 G_j(z)=\sum_{k\ge0}C_j(T_k)z^k\quad(j=3,4,5)
\end{equation}
and define
\begin{align}
 B(z)&=1-z-z^2-z^3,\notag\\
 Q(z)&=1-2z^3-2z^7-2z^{11}+z^{12},\label{eq:gfparts}\\
 D(z)&=(1-z)^2B(z)Q(z).\notag
\end{align}

\begin{proposition}[Exact rational cutoff series]\label{prop:GF}
For $j\in\{3,4,5\}$,
\begin{equation}\label{eq:GF}
 G_j(z)=\frac{P_j(z)}{D(z)},
\end{equation}
where the integer coefficients of $P_j$ are given in
Appendix~\ref{app:numerators}. Their degrees are $27,30,30$, respectively.
In addition,
\begin{equation}\label{eq:GFtotal}
 \frac1{1-z}+G_3(z)+G_4(z)+G_5(z)
   =\frac{1+z+z^2}{B(z)}.
\end{equation}
The common denominator displayed here is sufficient for the assertions;
no minimality claim is needed.
\end{proposition}
\begin{proof}[Finite-certificate proof]
For a proposed numerator $P_j$, compare its coefficients with those of
$D(z)\sum_{k\ge0}e_0A^kh_jz^k$. The verifier checks the coefficient
equalities through degree $106$. The degree of $D$ is $17$ and the largest
numerator degree is $30$.

For $k\ge17$, the residual coefficient is a fixed linear combination of
$e_0A^{k-i}h_j$ for $0\le i\le17$. Thus the residual sequence, after this
initial offset, satisfies the monic order-$76$ recurrence supplied by
$p(A)=0$. Its $76$ consecutive values at $k=31,\ldots,106$ are all zero.
The recurrence forces every subsequent value to be zero. The earlier
coefficients have been checked individually, proving \eqref{eq:GF}
identically, rather than to a finite truncation only.

The output $1$ occurs exactly for the zero representation, contributing
$1$ at every cutoff $T_k$. Summing all four output classes gives $T_k$.
Equation~\eqref{eq:Tgf} proves \eqref{eq:GFtotal}.
\end{proof}

The dominant-pole ratios provide an additional independent check on the
\emph{evaluation} of the three densities. Put $\rho=\beta^{-1}$ and
write $d_j=(a_j\beta^2+b_j\beta+c_j)/22$, with
\[
 (a_3,b_3,c_3)=(21,-38,5),\quad
 (a_4,b_4,c_4)=(-55,72,67),\quad
 (a_5,b_5,c_5)=(34,-34,-50).
\]
The verifier checks the exact polynomial divisibilities
\begin{equation}\label{eq:residuecert}
 B(z)\ \bigm|\ 22z^2P_j(z)
 -(a_j+b_jz+c_jz^2)(1-z)^2Q(z)(1+z+z^2).
\end{equation}
Evaluating at $\rho$ says precisely that the ratio of the dominant pole
coefficients in $G_j$ and the total cutoff generating function is $d_j$.
This agrees with the eigenvector aggregation. It is not used as a substitute
for the arbitrary-cutoff argument in Section~\ref{sec:allcutoffs}.

\section{Verification, reproducibility and limits of the claim}
\label{sec:verification}

\subsection{Finite proofs versus finite experiments}
The verifier contains two logically different types of check.
First, it verifies finite identities that support an all-index proof:
closure under both co-decomposition transitions; the full quotient diagram;
the right and left eigenvector equations; the integer annihilator; the
spectral separation inequality; and the generating-function recurrences.
Once verified, these combine with the induction and asymptotic arguments
in the article to establish claims for all lengths and all cutoffs.

Second, it performs additional finite cross-checks. These are useful for
detecting implementation mistakes but are not, by themselves, evidence
that replaces an all-index proof. The most substantial of these is a
separate enumeration of \emph{all} length-$n$ factors for every
$1\le n\le2048$.

\begin{lemma}[Complete finite coverage for the cross-check]\label{lem:paircheck}
The legal adjacent letter pairs of $\ttt$ are exactly
\[
 \mathcal P=\{00,01,02,10,20\}.
\]
If $m$ is chosen so that $\min_a|\varphi^m(a)|\ge H$, then every factor
of $\ttt$ of length at most $H$ occurs in at least one word
$\varphi^m(ab)$ with $ab\in\mathcal P$. Conversely, every factor of
one of those five words occurs in $\ttt$.
\end{lemma}
\begin{proof}
All five listed pairs occur in $\varphi^3(0)=0102010$. Inside a single
substitution image, the only pairs are $01$ and $02$. Across two images,
the next image begins with $0$ and the previous one ends in $0,1$, or $2$,
so the only possibilities are $00,10,20$. This proves the exact pair set.

Since $\ttt=\varphi^m(\ttt)$, the fixed point is a concatenation of blocks
$\varphi^m(a)$. A factor of length at most $H$ starting inside one block
cannot extend beyond the end of the next block: that next block already
has length at least $H$. Therefore the factor is contained in an image of
a legal adjacent pair. Conversely, each legal adjacent pair occurs in
$\ttt$, so its image occurs in $\varphi^m(\ttt)=\ttt$.
\end{proof}

For $H=2048$, the verifier uses $m=14$, for which the shortest image has
length $3136$. It forms prefix-sum arrays for the five pair images, slides
a window of each length through each image, and collects the exact sums.
This examines $93{,}070{,}336$ windows in total. The resulting number of
sums agrees with the automaton at every tested length. Unlike testing a
long guessed prefix, this enumeration has a proved coverage guarantee.

\subsection{Recorded verification results}
The standard run recorded in \code{data/verification\_results.json}
returns \code{PASS}. Its principal checks are:
\begin{center}
\small
\begin{tabular}{@{}p{0.61\textwidth}r@{}}
\toprule
Check & Size or outcome\\
\midrule
Co-decomposition pairs / closed sets & $56\,/\,277$\\
Reachable valid-language product states & $296$\\
Live quotient states / closed component & $76\,/\,66$\\
Left-eigenvector integer component equalities & $228$\\
Entries verified zero in $p(A)$ & $5776$\\
Independent complete factor-length checks & $2048$\\
Windows examined in the all-factor check & $93{,}070{,}336$\\
Individually evaluated lengths for count checks & $100{,}000$\\
Arbitrary prefix-count comparisons & $301$\\
Tribonacci cutoff comparisons & $61$\\
Generating-function coefficients through degree & $106$\\
Independent density-residue polynomial identities & $3$\\
\bottomrule
\end{tabular}
\end{center}
The matrix and algebraic checks use exact integers or rational numbers.
Decimal expansions are generated only for presentation and do not feed
back into a certificate.

\subsection{How to reproduce the package}
From the extracted package directory, run:
\begin{lstlisting}
python3 code/verify.py
python3 code/tribonacci.py --value 123456789
python3 code/tribonacci.py --count 1000000000000000000
\end{lstlisting}
The verifier needs no external Python package, no network access, and no
Walnut installation. It rebuilds the automaton and compares it with the
stored certificate instead of simply trusting the serialized table.
Its checks use explicit exceptions, not Python assertions, so they remain
active when invoked with \code{python3 -O}.

The optional script \code{code/derive\_spectrum.py} uses SymPy to reproduce
the symbolic derivation of the annihilator, the left eigenvector and the
rational generating functions. It is a discovery/reproduction aid, not a
dependency of the standard-library verification. It was tested with
SymPy~1.14.0. The script \code{code/make\_tables.py} regenerates the printed
table fragments and the exact-count data. The article can be built with
\begin{lstlisting}
pdflatex -halt-on-error article.tex
pdflatex -halt-on-error article.tex
\end{lstlisting}
A \code{Makefile} provides the same operations. The original random draw is
preserved rather than repeated during verification; its list digest and
selected index are checked for consistency.

\subsection{Scope, trust and remaining questions}
The claim is a complete proposed resolution of the stated natural-density
question, supported by an infinite argument and finite exact certificates.
It is not an independently refereed theorem, a Lean formalization, or a
proof that no previous resolution exists. The trusted computational layer
includes the Python interpreter and the executed verifier. The finite
tables and source code are provided so that another implementation can
check the same identities independently.

Several broader questions are not settled here: the optimal discrepancy
exponent; an optimal explicit error constant; the corresponding density
spectrum for arbitrary letter weights; and general density theorems for
all substitutive additive-complexity sequences. In particular, the general
regularity questions discussed in \cite{PSS,Couvreur} are not being
relabelled as consequences of this specific computation.

\section{Conclusion}
The additive-complexity proportions of the Tribonacci word are governed by
an algebraic left eigenvector of a finite transition matrix. The exact
weights are not uniform over recurrent states, but their output-class sums
collapse to three small expressions in the Tribonacci constant. A spectral
gap holds uniformly from every live state, and the lexicographic
prefix-cylinder decomposition transfers that estimate to every integer
cutoff. This proves the proposed densities and the power-saving estimate,
while the archived word-level and integer-matrix certificates make the
finite part of the argument reproducible.

\appendix
\clearpage
\section{The complete 76-state transition table}
\label{app:transitions}
The initial state is $0$. Read digits most significant first. The column
$o$ is the output, and a dash is an illegal transition to the omitted sink.
State $0$ has output $1$ and loops on $0$, accounting for leading zeros.
The two halves of the table together list every state exactly once.
\begin{center}
\small
\setlength{\tabcolsep}{9pt}
\renewcommand{\arraystretch}{1.10}
\begin{tabular}{@{}rrrr@{\hspace{14mm}}rrrr@{}}
\toprule
$q$ & $\delta(q,0)$ & $\delta(q,1)$ & $o$ &
$q$ & $\delta(q,0)$ & $\delta(q,1)$ & $o$\\
\midrule
\csname @@input\endcsname tables/transitions_rows.tex 
\bottomrule
\end{tabular}
\end{center}

\clearpage
\section{The exact normalized left eigenvector}
\label{app:eigenvector}
The table gives the nonzero entries of
$\ell_q=(u_q+w_q\beta+z_q\beta^2)/44$. The entries at
$0,1,2,3,4,6,7,8,12,13$ are zero. All coefficients are integers;
\eqref{eq:triplerule} is sufficient to verify $\ell A=\beta\ell$.
\begin{center}
\small
\setlength{\tabcolsep}{8pt}
\renewcommand{\arraystretch}{1.12}
\begin{tabular}{@{}rrrr@{\hspace{14mm}}rrrr@{}}
\toprule
$q$ & $u_q$ & $w_q$ & $z_q$ & $q$ & $u_q$ & $w_q$ & $z_q$\\
\midrule
\csname @@input\endcsname tables/eigenvector_rows.tex 
\bottomrule
\end{tabular}
\end{center}
\noindent
The coefficient sums are $(44,0,0)$, so $\ell\one=1$ exactly. Summing only
rows with a specified output gives the three density formulas in
Theorem~\ref{thm:main}. The machine-readable version is
\code{data/eigenvector.json}.

\clearpage
\section{Numerators of the rational cutoff generating functions}
\label{app:numerators}
For $j=3,4,5$, write $P_j(z)=\sum_{k=0}^{30}p_{j,k}z^k$ and use the common
denominator $D(z)$ of \eqref{eq:gfparts}. Coefficients beyond the displayed
range are zero. This table, together with $D$, fully specifies
Proposition~\ref{prop:GF}.
\begin{center}
\small
\setlength{\tabcolsep}{14pt}
\renewcommand{\arraystretch}{1.10}
\begin{tabular}{@{}rrrr@{}}
\toprule
$k$ & $p_{3,k}$ & $p_{4,k}$ & $p_{5,k}$\\
\midrule
\csname @@input\endcsname tables/numerator_rows.tex 
\bottomrule
\end{tabular}
\end{center}

\clearpage
\section{Artifact map and selection record}
\label{app:artifacts}
The core package has the following organization. The README gives individual
filenames and command-line examples for every script.
\begin{center}
\small
\begin{tabular}{@{}p{0.36\textwidth}p{0.59\textwidth}@{}}
\toprule
Path & Purpose\\
\midrule
\code{article.tex}, \code{article.pdf} & Article source and typeset report.\\
\code{code/tribonacci.py} & Exact reconstruction, single-value evaluation and prefix counting.\\
\code{code/verify.py} & Standard-library verifier of the finite certificates and cross-checks.\\
Other scripts in \code{code/} & Optional symbolic derivation, table regeneration and random-selection auditing.\\
Automaton data in \code{data/} & All $277$ co-decomposition sets, the $76$-state machine and its quotient information.\\
Algebraic data in \code{data/} & Exact eigenvector triples, generating functions and decimal constants.\\
Count data in \code{data/} & Exact counts at large cutoffs, including $10^{100}$, in JSON and CSV form.\\
Verification logs in \code{data/} & Recorded outcomes, bounds and actual verification scope.\\
\code{data/selection.json} & Original ordered $120$-area list and the single recorded selection.\\
\code{notes/} & Source-version and manifest-exclusion audits, plus the original area-draw script.\\
\code{STATUS.md}, \code{README.md} & Claim boundaries and practical reproduction instructions.\\
\bottomrule
\end{tabular}
\end{center}

The original draw selected item $68$, \emph{Combinatorics on words}.
The SHA-256 of the ordered list, encoded as UTF-8 strings separated by a
single newline, is
\begin{center}\small\ttfamily
 ea66507989230ed06b68907eef7e2e8c14f\par
 933ba70aa97b56a70fab0ea4477b6.
\end{center}
The two lines above are one concatenated digest, with no intervening space.
The digest authenticates the recorded list against accidental alteration;
it is not an external timestamp or a cryptographic proof of the historical
random-draw procedure. The record reports the actual executed single draw.

\begin{thebibliography}{9}
\bibitem{PSS}
Pierre Popoli, Jeffrey Shallit and Manon Stipulanti,
\emph{Additive Word Complexity and Walnut}.
In \emph{44th IARCS Annual Conference on Foundations of Software Technology
and Theoretical Computer Science (FSTTCS 2024)},
LIPIcs \textbf{323} (2024), Article~32, 32:1--32:18.
\href{https://doi.org/10.4230/LIPIcs.FSTTCS.2024.32}{doi:10.4230/LIPIcs.FSTTCS.2024.32}.
Full preprint: \href{https://arxiv.org/abs/2410.02409v1}{arXiv:2410.02409v1},
3 October 2024. The theorem and remark numbers in this article refer to
Theorem~18 and Remark~19 of that full preprint.

\bibitem{Turek}
Ond\v{r}ej Turek,
\emph{Abelian Complexity Function of the Tribonacci Word}.
Journal of Integer Sequences \textbf{18} (2015), Article~15.3.4.
\url{https://cs.uwaterloo.ca/journals/JIS/VOL18/Turek/turek3.html}.
Also \href{https://arxiv.org/abs/1309.4810v2}{arXiv:1309.4810v2}.
The co-decomposition approach is prior work; its specialization needed here
is proved in Section~\ref{sec:automaton}.

\bibitem{Couvreur}
Jean-Michel Couvreur, Martin Delacourt, Nicolas Ollinger, Pierre Popoli,
Jeffrey Shallit and Manon Stipulanti,
\emph{Effective Computation of Generalized Abelian Complexity for Pisot Type
Substitutive Sequences}.
\href{https://arxiv.org/abs/2504.13584v2}{arXiv:2504.13584v2},
22 April 2025. Consulted for the current literature audit and related
automata constructions; not imported as a density theorem.
\end{thebibliography}

\end{document}
